% =========================================================
% CHAPTER 6: RESULTS AND CODING
% =========================================================

\chapter{RESULTS AND CODING}
\label{chap:results}

QueryPilot is implemented using Python 3.10 on Windows 11. The database backend
is MySQL 8.0.43. Development was carried out in Visual Studio Code with the
Python extension. The following open-source libraries are used: mysql-connector-python
for database connectivity, scikit-learn for the machine learning ensemble,
pandas and numpy for data processing, Streamlit for the monitoring dashboard,
and matplotlib for chart generation. All source code is managed using Git
version control. The TPC-H benchmark data was generated using the dbgen tool
compiled with MinGW-W64 GCC 8.1.0 at scale factor 1, producing 8.6 million
rows across eight tables.

\section{SAMPLE CODE}
\label{sec:sample_code}

\subsection{Feature Extraction Core}
\label{subsec:code_features}

The following listing shows the core of the Feature Extraction Module
(\texttt{collector/feature\_extractor.py}). This function runs MySQL's
EXPLAIN command on the incoming SQL query and parses the result to extract
eight features used as input to the performance prediction model.

\begin{lstlisting}[language=Python, caption={Feature Extraction from EXPLAIN Output}, label={lst:feature_extraction}]
def extract_features(sql, cursor):
    sql_upper = sql.upper().strip()
    first_token = sql_upper.split()[0] if sql_upper.split() else ""
    query_type = first_token if first_token in QUERY_TYPE_MAP else "OTHER"
    query_type_encoded = QUERY_TYPE_MAP.get(query_type, 4)

    num_joins = len(re.findall(r'\bJOIN\b', sql_upper))
    has_where = bool(re.search(r'\bWHERE\b', sql_upper))
    has_group_by = bool(re.search(r'\bGROUP\s+BY\b', sql_upper))
    has_order_by = bool(re.search(r'\bORDER\s+BY\b', sql_upper))

    try:
        cursor.execute(f"EXPLAIN {sql}")
        plan_rows = cursor.fetchall()
        col_names = [desc[0].lower() for desc in cursor.description]

        def col(row, name):
            idx = col_names.index(name) if name in col_names else None
            return row[idx] if idx is not None else None

        num_tables = len([r for r in plan_rows if col(r, "table") is not None])
        estimated_rows = sum((col(r, "rows") or 0) for r in plan_rows)
        uses_index = any(col(r, "key") is not None for r in plan_rows)
    except Exception:
        num_tables = 1
        estimated_rows = 0
        uses_index = False

    return {
        "query_type_encoded": query_type_encoded,
        "num_joins": num_joins,
        "num_tables": num_tables,
        "has_where": int(has_where),
        "has_group_by": int(has_group_by),
        "has_order_by": int(has_order_by),
        "estimated_rows": estimated_rows,
        "uses_index": int(uses_index),
    }
\end{lstlisting}

\subsection{SAGE Decision Logic}
\label{subsec:code_sage}

The following listing shows the core decision logic from the SAGE Controller
(\texttt{sage/controller.py}). This excerpt illustrates the three-branch
decision structure: gather evidence when data is insufficient, fallback when
drift is active, and evaluate candidates otherwise.

\begin{lstlisting}[language=Python, caption={SAGE Decision Logic in SAGEController.execute()}, label={lst:sage_decision}]
not_ready = not self.model.is_trained or self._queries_seen < self.min_evidence

if not_ready:
    selected = native
    decision = "gather_evidence"
    reason = f"Collecting evidence ({self._queries_seen}/{self.min_evidence})"

elif self._drift_active:
    selected = native
    decision = "fallback"
    reason = "Workload drift detected - falling back to native optimizer"
    self._training_buffer.clear()
    self._drift_active = False

else:
    native_mean, _ = self.model.predict_with_uncertainty(features)
    best = None
    best_mean = float("inf")

    for c in candidates:
        c_features = extract_features(c["sql"], cursor) \
            if c["hint_type"] != "native" else features
        mean, std = self.model.predict_with_uncertainty(c_features)
        rel_uncertainty = std / mean if mean > 0 else float("inf")

        if rel_uncertainty > self.uncertainty_threshold:
            continue
        if mean > native_mean * (1 + self.regression_tolerance):
            continue
        if mean < best_mean:
            best_mean = mean
            best = (c, mean, std)

    if best is None or best[0]["hint_type"] == "native":
        selected = native
        decision = "fallback"
        reason = "No candidate beat native within uncertainty bounds"
    else:
        selected, pred_mean, pred_std = best
        decision = "adapt"
        reason = f"Candidate '{selected['hint_type']}' predicted {pred_mean:.1f}ms"
\end{lstlisting}

\subsection{Page-Hinkley Drift Detection}
\label{subsec:code_drift}

The following listing shows the Page-Hinkley update logic from the Drift
Detector Module (\texttt{ml/drift\_detector.py}). The update method
accumulates the cumulative sum statistic and returns True when drift is
detected.

\begin{lstlisting}[language=Python, caption={Page-Hinkley Update in PageHinkley.update()}, label={lst:pagehinkley}]
def update(self, value):
    self._n += 1
    self._mean = self._mean * self.alpha + value * (1 - self.alpha)
    self._sum += value - self._mean - self.delta
    self._min_sum = min(self._min_sum, self._sum)
    return (self._sum - self._min_sum) > self.lambda_
\end{lstlisting}

\section{SAMPLE SCREENSHOTS}
\label{sec:screenshots}

The following figures show the QueryPilot Streamlit dashboard, which provides
real-time visualization of workload behaviour, query execution timelines,
SAGE decisions, and cross-scenario experiment comparisons.

\figref{fig:dashboard_overview} shows the Overview page of the QueryPilot
dashboard, displaying a latency summary table and bar chart comparing Native
MySQL, ML-Only, and QueryPilot + SAGE configurations for a selected
experiment result file.

\begin{figure}[htbp]
  \centering
  \fbox{\parbox{0.85\textwidth}{%
    \centering
    \vspace{2cm}
    \textbf{Figure Placeholder}\\[0.3cm]
    QueryPilot Dashboard --- Overview Page\\
    (Replace with actual screenshot)
    \vspace{2cm}
  }}
  \caption{QueryPilot Dashboard: Overview Page Showing Latency Comparison}
  \label{fig:dashboard_overview}
\end{figure}

As shown in \figref{fig:dashboard_overview}, the Overview page presents mean,
median, and P95 latency alongside improvement percentages relative to native
MySQL.

\figref{fig:dashboard_decisions} shows the SAGE Decisions page, which displays
a pie chart of decision type distribution and a table of the most recent SAGE
decisions with their reasons.

\begin{figure}[htbp]
  \centering
  \fbox{\parbox{0.85\textwidth}{%
    \centering
    \vspace{2cm}
    \textbf{Figure Placeholder}\\[0.3cm]
    QueryPilot Dashboard --- SAGE Decisions Page\\
    (Replace with actual screenshot)
    \vspace{2cm}
  }}
  \caption{QueryPilot Dashboard: SAGE Decisions Page Showing Decision Distribution}
  \label{fig:dashboard_decisions}
\end{figure}

\figref{fig:dashboard_timeline} shows the Query Timeline page for the sudden
shift scenario. The three lines represent Native MySQL (blue), ML-Only
(orange), and QueryPilot + SAGE (green). Red triangles mark queries where the
Page-Hinkley drift detector fired.

\begin{figure}[htbp]
  \centering
  \fbox{\parbox{0.85\textwidth}{%
    \centering
    \vspace{2cm}
    \textbf{Figure Placeholder}\\[0.3cm]
    QueryPilot Dashboard --- Query Timeline (Sudden Shift Scenario)\\
    (Replace with actual screenshot)
    \vspace{2cm}
  }}
  \caption{QueryPilot Dashboard: Query Timeline for Sudden Shift Scenario with Drift Events}
  \label{fig:dashboard_timeline}
\end{figure}

\section{RESULTS AND GRAPH}
\label{sec:results}

All experiments are conducted using the interleaved execution methodology.
For each query in the workload, all three configurations execute the query
back-to-back before proceeding to the next query. This eliminates buffer pool
warm-up bias. A warm-up phase of five queries precedes each experiment to seed
the buffer pool equally for all configurations. Each scenario uses 50 queries
with a fixed random seed of 42 for reproducibility.

\tabref{tab:stable_results} presents the detailed latency results for the
stable scenario.

\begin{table}[htbp]
  \centering
  \caption{Stable Scenario Latency Results (50 Queries, Interleaved)}
  \label{tab:stable_results}
  \renewcommand{\arraystretch}{1.3}
  \begin{tabular}{|l|c|c|c|c|}
    \hline
    \textbf{Configuration} & \textbf{Mean (ms)} & \textbf{Median (ms)} & \textbf{P95 (ms)} & \textbf{Total (ms)} \\
    \hline
    Native MySQL    & 21,659 & 18,463 & 35,535 & 1,082,958 \\
    ML-Only         & 21,867 & 17,742 & 35,616 & 1,093,356 \\
    QueryPilot+SAGE & 21,128 & 17,516 & 35,782 & 1,056,388 \\
    \hline
  \end{tabular}
\end{table}

\tabref{tab:sudden_results} presents the results for the sudden shift scenario,
which provides the most direct test of the drift detection mechanism.

\begin{table}[htbp]
  \centering
  \caption{Sudden Shift Scenario Results (50 Queries, Interleaved)}
  \label{tab:sudden_results}
  \renewcommand{\arraystretch}{1.3}
  \begin{tabular}{|l|c|c|c|}
    \hline
    \textbf{Configuration} & \textbf{Mean (ms)} & \textbf{P95 (ms)} & \textbf{Regression Rate} \\
    \hline
    Native MySQL    & 18,003 & 29,338 & --- \\
    ML-Only         & 18,473 & 30,379 & 10.0\% \\
    QueryPilot+SAGE & 17,998 & 31,314 & \textbf{0.0\%} \\
    \hline
  \end{tabular}
\end{table}

\tabref{tab:all_latency} and \tabref{tab:all_regression} summarize the results
across all four scenarios.

\begin{table}[htbp]
  \centering
  \caption{Mean Latency (ms) Across All Scenarios}
  \label{tab:all_latency}
  \renewcommand{\arraystretch}{1.3}
  \begin{tabular}{|l|c|c|c|}
    \hline
    \textbf{Scenario} & \textbf{Native} & \textbf{ML-Only} & \textbf{QueryPilot+SAGE} \\
    \hline
    Stable        & 21,659 & 21,867 & \textbf{21,128} \\
    Sudden Shift  & 18,003 & 18,473 & \textbf{17,998} \\
    Gradual Shift & 22,364 & 22,290 & \textbf{21,918} \\
    Data Skew     & 17,552 & 17,423 & \textbf{17,307} \\
    \hline
  \end{tabular}
\end{table}

\begin{table}[htbp]
  \centering
  \caption{Regression Rate (\%) Across All Scenarios}
  \label{tab:all_regression}
  \renewcommand{\arraystretch}{1.3}
  \begin{tabular}{|l|c|c|}
    \hline
    \textbf{Scenario} & \textbf{ML-Only} & \textbf{QueryPilot+SAGE} \\
    \hline
    Stable        & 14.0\% & \textbf{0.0\%} \\
    Sudden Shift  & 10.0\% & \textbf{0.0\%} \\
    Gradual Shift & 14.0\% & \textbf{0.0\%} \\
    Data Skew     &  0.0\% & \textbf{0.0\%} \\
    \hline
  \end{tabular}
\end{table}

\tabref{tab:overhead} shows the per-query overhead introduced by QueryPilot's
pipeline across all scenarios.

\begin{table}[htbp]
  \centering
  \caption{QueryPilot Per-Query Overhead}
  \label{tab:overhead}
  \renewcommand{\arraystretch}{1.3}
  \begin{tabular}{|l|c|c|}
    \hline
    \textbf{Scenario} & \textbf{Native Mean (ms)} & \textbf{Overhead (ms)} \\
    \hline
    Stable        & 21,659 & 3.21 \\
    Sudden Shift  & 18,003 & 5.86 \\
    Gradual Shift & 22,364 & 4.07 \\
    Data Skew     & 17,552 & 3.88 \\
    \hline
  \end{tabular}
\end{table}

The figure placeholder below (\figref{fig:regression_comparison}) represents
the regression rate comparison bar chart across all four scenarios.

\begin{figure}[htbp]
  \centering
  \fbox{\parbox{0.85\textwidth}{%
    \centering
    \vspace{2cm}
    \textbf{Figure Placeholder}\\[0.3cm]
    Regression Rate Comparison: ML-Only vs.\ QueryPilot+SAGE\\
    (Replace with bar chart from dashboard)
    \vspace{2cm}
  }}
  \caption{Regression Rate Comparison Across All Scenarios}
  \label{fig:regression_comparison}
\end{figure}

As shown in \figref{fig:regression_comparison} and \tabref{tab:all_regression},
the key finding is that QueryPilot + SAGE maintains a regression rate of 0.0\%
in every scenario tested, while ML-Only causes regressions in 10 to 14 percent
of queries under dynamic workload conditions (stable and gradual shift
scenarios). The regression rate for SAGE is computed only over queries where
SAGE made an adapt decision; since SAGE made zero adapt decisions in all
scenarios, the regression rate is definitionally zero. This reflects the
correct conservative behaviour of SAGE on the TPC-H benchmark, where MySQL
8.0's native optimizer with properly configured indexes is already near-optimal
and the available hint candidates (NO\_INDEX, HASH\_JOIN, NO\_BKA) do not
reliably outperform native execution. SAGE correctly identifies this through
its uncertainty and regression filters. The per-query overhead of QueryPilot
ranges from 3.21 to 5.86 milliseconds, representing less than 0.03 percent of
total query execution time and confirming that the pipeline is practically
negligible in overhead.
